\subsection{The universal modular comparison and its precise remaining lifts}

\begin{statement}{B.16}{Open}{Central twists, normalized conjugators, and the universal source comparison}
In this subsection use the normalization of C3:572--582,
\[
 C=\widehat{\SL_2(\Z)},\quad
 r=\begin{pmatrix}0&1\\-1&1\end{pmatrix},\quad
 s=\begin{pmatrix}0&-1\\1&0\end{pmatrix},\quad
 \omega=r^3=s^2=-I,\quad
 e=sr=\begin{pmatrix}1&-1\\0&1\end{pmatrix}.
 \tag{B.16.1}
\]
Thus $s$ is the negative of the final \S48 matrix $\sigma$.
All groups and automorphisms in this subsection are profinite and
continuous; they are not integral group schemes. Put
\[
 \overline C=C/\langle\omega\rangle,
 \quad P=\ker(C\longrightarrow\SL_2(\F_2)),\quad
 L=\overline{\langle e\rangle},\quad
 P_0=\overline{\langle\omega e^2,\omega(rs)^2\rangle}.
 \tag{B.16.2}
\]
The following assertions are proved below, using the explicitly
identified torsion and centralizer theorems of Herfort--Ribes:
\[
\begin{gathered}
 Z(C)=\langle\omega\rangle,
 \quad C^{\rm ab}=\Z/12\quad(r\mapsto2, s\mapsto3),\qquad
 L\simeq\widehat\Z,\quad P=P_0\times\langle\omega\rangle,\\
 1\longrightarrow\langle\Theta\rangle\longrightarrow\Aut(C)
    \longrightarrow\Aut(\overline C)\longrightarrow1,
 \qquad \Theta(x)=x\omega^{\operatorname{sg}(x)},
 \tag{B.16.3}
\end{gathered}
\]
where $\operatorname{sg}:C\to\Z/2$ is the sign after reduction to
$\SL_2(\F_2)=S_3$. In particular $\Theta(r)=r$ and
$\Theta(s)=\omega s$. It fixes $L^2$ pointwise, but does not carry
$L$ to a $C$-conjugate of $L$.

Let $\mathcal A$ consist of the automorphisms of $C$ preserving $P$
and the conjugacy class of $L$, and let $\overline{\mathcal A}$ be
the analogous subgroup of $\Aut(\overline C)$ for $P/\omega$ and
the image of $L$. Then
\[
                  \mathcal A\xrightarrow{\sim}\overline{\mathcal A}.
 \tag{B.16.4}
\]
Let $Z$ be the group of automorphisms preserving $L$ and the two
conjugacy classes of cyclic subgroups $\langle r\rangle,
\langle s\rangle$. The multiplier $\mu$ and the two signs are
unique and satisfy
\[
 u(e)=e^\mu,\quad u(r)\sim_C r^{\epsilon_3(\mu)},\quad
 u(s)\sim_C s^{\epsilon_2(\mu)}.
 \tag{B.16.5}
\]
Consequently the independent-sign and linked-sign groups of B.13
coincide in this particular universal case. Define
\[
 \gamma=Z/\operatorname{Int}(L),\qquad
 G=(C\rtimes Z)/\{(\ell^{-1},\operatorname{Int}\ell):\ell\in L\}.
 \tag{B.16.6}
\]
The adjoint action induces a topological isomorphism
$G/\langle\omega\rangle\simeq\mathcal A$.

Both $P$ and $P_0$ are normal in $G$. The action on $C/P=S_3$
defines $\kappa:G\to\Aut(S_3)=S_3$; put $G^!=\ker\kappa$ and
$Z^!=Z\cap G^!$. There are exact sequences
\[
\begin{aligned}
 1&\longrightarrow L^2\longrightarrow Z^!
       \longrightarrow\gamma\longrightarrow1,\\
 1&\longrightarrow P\longrightarrow G^!
       \longrightarrow\gamma\longrightarrow1.
\end{aligned}
 \tag{B.16.7}
\]
In particular $Z^!$ has index two in $Z$; the distinction corrected
in the margin of C3:582 and again at C3:613 is necessary.

Write $\tau$ for conjugation by $\operatorname{diag}(-1,1)$ and
$\tau'=s^{-1}\tau$ in $G$, as in B.13. For every $u\in Z^!$ there
are unique
\[
 \alpha\in P_0\tau^{\nu_3(\mu)},\qquad\beta\in P_0,
 \qquad (-1)^{\nu_i(\mu)}=\epsilon_i(\mu),
 \tag{B.16.8}
\]
such that
\[
 u(r)=\alpha r\alpha^{-1},\qquad
 u(s)=\omega^{\nu_2(\mu)}\beta s\beta^{-1}.
 \tag{B.16.9}
\]
This proves the normalization of C3:579(27),582(38)--(40) for
the precisely defined subgroup $Z^!$.

Identify $N_r^*,N_s^*$ with the subgroups of $G$ in B.13. The maps
of C3:612(118)--(123) have the following rigorous formal version:
\[
 \overline a(u)=\tau'^{\nu_3(\mu)}\alpha^{-1}u\pmod\omega,
 \qquad b(u)=\beta^{-1}u,
 \qquad
 \overline a:Z^!\to N_r^*/\omega,\quad b:Z^!\to N_s^*.
 \tag{B.16.10}
\]
They are continuous homomorphisms over $\gamma$ and satisfy
\[
 \overline a(e^{2t})=1,\qquad b(e^{2t})=\omega^{t\bmod2}
 \quad(t\in\widehat\Z).
 \tag{B.16.11}
\]
Here $e^{2t}$ in the domain means its embedded inner automorphism.
They induce sections
\[
 a_\gamma:\gamma\to N_r^*/\omega,\qquad
 b_\gamma:\gamma\to N_s^*/\omega,
 \tag{B.16.12}
\]
and therefore
\[
 N_r^*/\omega\simeq C_3\rtimes_{\epsilon_3}\gamma,
 \qquad N_s^*/\omega\simeq C_2\times\gamma.
 \tag{B.16.13}
\]
These are the quotient sections in C3:615(130)--(131). The stronger
unquotiented $\sigma$-section of (132)--(140), its chosen geometric
interpretation, and the source comparisons specified below are not
claimed by this statement.
\end{statement}

\begin{proof}[Proved part]
We first specify the sole external group-theoretic input.
Herfort and Ribes, \emph{On Free Pro-$\mathcal C$-Products of
Pro-$\mathcal C$-Groups}, Theorems 1 and 2, pp.~31--34
\cite{HerfortRibes1984}, apply to the class of all finite groups.
For a free profinite product they put finite-order elements in
conjugates of the factors and give $C_{A\amalg B}(a)=C_A(a)$ for
$1\ne a\in A$. Their hypotheses require closure under subgroups,
quotients and extensions, all of which the class of all finite
groups satisfies. The directly inspected primary article is an
announcement with a proof outline; it explicitly refers to their
subsequent full paper. We use these two established theorems,
not a claim that a discrete reduced-word proof works profinitely.

By C.2 and the universal property of completion,
$\overline C=C_3\amalg C_2$. Indeed the completion of the discrete
quotient by its finite central subgroup is $C/\omega$: both groups
have exactly the same continuous maps to finite groups. The
involution survives modulo three, so the kernel has order two.
The centralizer theorem, first for the order-three factor and then
for the order-two factor, gives
\[
 C_C(r)=\langle r\rangle,\qquad
 C_C(s)=\langle s\rangle,\qquad Z(C)=\langle\omega\rangle.
 \tag{B.16.14}
\]
For example an element centralizing $r$ projects into
$\langle\overline r\rangle$, whose inverse image is exactly
$\langle r\rangle$; the reverse inclusion is immediate.

Abelianizing the presentation in C.2 gives
$\langle r,s\mid6r=4s=0, 3r=2s\rangle$.
Its order is twelve, and the displayed map to $\Z/12$ is onto,
so it is an isomorphism. The profinite abelianization has the same
universal property for finite abelian targets, hence is that same
finite group. If $v$ fixes $\omega$ and induces the identity on
$\overline C$, write $v(x)=x\omega^{\chi(x)}$. Centrality makes
$\chi:C\to\Z/2$ a continuous homomorphism; fixing $\omega$
requires $\chi(\omega)=0$. Conversely such a character defines
an involutory automorphism. There are exactly two characters,
the trivial one and $\operatorname{sg}$, and both vanish on
$\omega$, whose abelianization is $6$. This proves the kernel
in (B.16.3), and also proves the same assertion for endomorphisms
inducing the identity on the central quotient.

For surjectivity in (B.16.3), let $v\in\Aut(\overline C)$.
The torsion theorem gives conjugators for
$v(\overline r)\sim\overline r^{\pm1}$ and
$v(\overline s)\sim\overline s$. Lift those conjugators to $C$.
The resulting proposed images of $r,s$ have common cube and square
$\omega$, so C.2 and completion extend them to an endomorphism
$\widetilde v$ of $C$. Make the same construction for $v^{-1}$.
Both compositions induce the identity on $\overline C$ and fix
$\omega$; the preceding kernel calculation makes them
automorphisms. Thus $\widetilde v$ is an automorphism too.

The element $e$ has order $n$ in every $\SL_2(\Z/n)$.
Consequently $t\mapsto e^t$ embeds $\widehat\Z$ into $C$.
Modulo three, $e$ has order three and $\omega e$ has order six.
Hence $\Theta(L)$ and $L$ cannot be conjugate, whereas
$\Theta(e^2)=e^2$. An automorphism of $\overline C$ preserving
the indicated peripheral subgroup class sends $\overline e$ to a
conjugate of $\overline e^{\mu}$ for a unique unit multiplier
after its conjugator is chosen. Either lift to $C$ sends $e$ to
that lift or its negative; composing with $\Theta$ switches these
two choices. Exactly one lift preserves the class of $L$.
Preservation of $P$ is equivalent to preservation of $P/\omega$.
This proves (B.16.4).

The two signs in (B.16.5) are unique, since $r,r^{-1}$ have
distinct images $2,10$ in $C^{\rm ab}$, and $s,s^{-1}$ have
distinct images $3,9$. In that abelianization the loop equation is
\[
                 2e_r+3e_s=5\mu\pmod {12}.
\]
Reducing modulo three gives $e_r\equiv\mu\pmod3$; reducing
modulo four gives $e_s\equiv\mu\pmod4$. This proves the two
linked signs. It also proves that the numerical coordinates of
B.13 introduce no kernel in this particular universal case.

E.15 proves that the discrete $\Pi_0$ is normal of index twelve,
free on $\lambda_0=\omega e^2,\lambda_1=\omega(rs)^2$, and
that $\Pi=\Pi_0\times\langle\omega\rangle$. A subgroup of finite
index carries its full profinite topology: the core in the ambient
group of any finite-index subgroup is still finite-index.
Completion therefore gives $P=P_0\times\langle\omega\rangle$,
with $P_0$ free profinite on $\lambda_0,\lambda_1$.
Equivalently, $P_0$ is the closed normal subgroup of $C$ generated
by $\lambda_0$; conjugation by $r$ supplies $\lambda_1$.
Every $u\in Z$ sends $\lambda_0$ to $\lambda_0^\mu$ and fixes
$\omega$. It therefore preserves both $P_0$ and $P$.

For completeness, $N_C(L)=L\langle\omega\rangle$.
An element of this normalizer has reduction modulo two in
$C_{S_3}(\overline e)=\langle\overline e\rangle$.
Multiplying by $e$ if necessary puts it in $P$. If an element of
$P$ carries $e^2$ to $e^{2v}$, its conjugation on
$P_0^{\rm ab}=\widehat\Z^2$ forces $v=1$.
The centralizer theorem applied to
$P_0=\langle\lambda_0\rangle\amalg\langle\lambda_1\rangle$
then puts it in $\langle\lambda_0\rangle\langle\omega\rangle$.
This is $L^2\langle\omega\rangle$, proving the claim.
In particular $L\cap Z(C)=1$.

The torsion theorem also says that every automorphism of $C$
carries $r$ to a conjugate of $r^{\pm1}$ and $s$ to a conjugate
of $s^{\pm1}$. For $r$, the possible extra central sign after
lifting from $\overline C$ is ruled out by its cube being
$\omega$. Thus $\mathcal A=\operatorname{Int}(C)Z$.
The action of $C\rtimes Z$ on $C$ has kernel consisting of
$(c,\operatorname{Int}(c^{-1}))$ with $c\in N_C(L)$.
The preceding normalizer computation and the antidiagonal quotient
in (B.16.6) leave exactly $\langle\omega\rangle$ as kernel.
All groups are compact Hausdorff, so the resulting continuous
bijection $G/\omega\to\mathcal A$ is an isomorphism.

Since $Z$ and $C$ preserve $P_0,P$, so does $G$. The action of
$u\in Z$ on $S_3$ fixes the transposition $\overline e$, because
$\mu$ is odd. Since $\Aut(S_3)=\operatorname{Int}(S_3)$, its
image under $\kappa$ lies in $\langle\overline e\rangle$.
The element $\operatorname{Int}(e)$ maps onto that order-two
subgroup. Thus $Z^!$ has index two, and multiplying by a suitable
element of $L$ makes any coset in $\gamma$ lie in $Z^!$.
Its kernel is $L^2$. Similarly $\kappa(C)=S_3$ corrects any
element of $G$ into $G^!$, whose intersection with $C$ is $P$.
This proves (B.16.7).

We now prove normalization without supposing it in advance.
For $u\in Z^!$, write $u(s)=c s^{\epsilon_2}c^{-1}$.
Modulo $P$ the conjugator $c$ centralizes $\overline s$;
multiply it by a suitable power of $s$ to put it in $P$.
Removing its central $\omega$-factor puts it in $P_0$, giving
$\beta$. For $r$ the same argument works when $\epsilon_3=1$,
using $C_{S_3}(\overline r)=\langle\overline r\rangle$.
When $\epsilon_3=-1$, use
$\tau(r)=sr^{-1}s^{-1}$. If $u(r)=cr^{-1}c^{-1}$,
then $cs^{-1}$ conjugates $\tau(r)$ to $u(r)$ and has reduction
in $\langle\overline r\rangle$. Multiplication by a suitable
power of $\tau(r)$, and then removal of the central sign,
produces $\alpha\in P_0\tau$.
Two such $\beta$ differ by $C_C(s)\cap P_0=1$; two such
$\alpha$ of the prescribed parity differ by
$C_C(r)\cap P_0=1$, by (B.16.14) and the injections of the
order-four and order-six factors into $C/P_0$ proved in E.15.
The coordinate maps are continuous: their graphs are closed,
and a continuous bijection from a compact space to a Hausdorff
space has continuous inverse.

The element $b(u)$ in (B.16.10) is characterized uniquely by
\[
             b(u)\in N_s^*,\qquad b(u)\equiv u\pmod {P_0}.
 \tag{B.16.15}
\]
Indeed $N_s^*\cap P_0=C_C(s)\cap P_0=1$.
Since $P_0$ is normal in $G$, (B.16.15) proves multiplicativity.
Likewise $a(u)=\tau'^{\nu_3}\alpha^{-1}u$ belongs to $N_r^*$,
lies over the same element of $\gamma$, and satisfies
\[
                   a(u)\equiv s^{-\nu_3}u\pmod P.
 \tag{B.16.16}
\]
Here $N_r^*\cap P=\langle\omega\rangle$, so the class of
$a(u)$ modulo $\omega$ is uniquely characterized by this
congruence. Elements of $Z^!$ commute with $C/P$ in $G/P$.
For $u',u\in Z^!$ their right sides in (B.16.16) multiply to
$s^{-(\nu_3'+\nu_3)}u'u$; replacing the sum by its value modulo
two changes a power of $s^2=\omega\in P$ only.
This proves that $\overline a$ is a homomorphism.

For $u=e^{2t}=\omega^{t\bmod2}\lambda_0^t$ the two
characterizations give (B.16.11). Hence the quotient maps
(B.16.12) exist and are sections over $\gamma$.
Their kernels in $N_r^*/\omega,N_s^*/\omega$ are respectively
$\langle r\rangle/\omega=C_3$ and
$\langle s\rangle/\omega=C_2$. Conjugation in the first case
is the sign $\epsilon_3$; every automorphism of $C_2$ is the
identity. This proves (B.16.13). In particular the reference to
$\epsilon_2$ acting on the order-two quotient in C3:615 must
be read as a trivial action; inversion becomes visible on $C_4$.
\end{proof}

\paragraph{Executed finite model and its exact limitation.}
The stand-alone program \texttt{AB\_modular\_comparison.py}
enumerates every matrix and every automorphism from the images of
$r,s$, checking every edge of the complete Cayley graph.
For $H_4=\SL_2(\Z/4)$ it gives
\[
 |H_4|=48,\quad H_4/\omega\simeq S_4,\quad
 |\Aut(H_4)|=48,\quad
 |\ker(\Aut(H_4)\to\Aut(S_4))|=2,
 \tag{B.16.17}
\]
and the image has order twenty-four. The isomorphism to $S_4$
is certified by conjugation on the four cyclic order-three
subgroups: the permutation image has order twenty-four and
kernel $\{1,\omega\}$. This also gives $|\Aut(S_4)|=24$:
an automorphism acts on those four intrinsic subgroups, and
an automorphism acting trivially there must commute with the
faithful inner permutation action, which forces it to fix every
group element. The central-character argument above bounds the
remaining kernel by two, while the sign twist realizes both.

In this finite group all forty-eight automorphisms preserve the
conjugacy class of $\langle e\rangle$. In fact
\[
 \begin{pmatrix}1&0\\2&1\end{pmatrix}
 e\begin{pmatrix}1&0\\2&1\end{pmatrix}^{-1}
       =-e^{-1}\pmod4.
 \tag{B.16.18}
\]
The finite counterpart of $P_0$ is the normal Klein four group
\[
 \left\{I,
 \begin{pmatrix}1&2\\2&1\end{pmatrix},
 \begin{pmatrix}3&2\\0&3\end{pmatrix},
 \begin{pmatrix}3&0\\2&3\end{pmatrix}\right\}.
\]
Using the numerical group $(\Z/12)^\times$, the program finds
sixteen linked automorphism--multiplier pairs. Eight act
trivially modulo two; each has one normalized $\alpha$ of the
prescribed parity and two choices of $\beta$. For the identity
with multiplier one the latter are already
$I$ and $\left(\begin{smallmatrix}1&2\\2&1\end{smallmatrix}\right)$.
Thus finite quotient centralizers do not prove profinite
normalization; the centralizer theorem is essential to our
reduction. The additional modulo-three computation confirms the
orders three and six used above. It also verifies that the
proposed central twist does not descend to an automorphism of
$\SL_2(\F_3)$; its central-quotient automorphism kernel is trivial.

\paragraph{Exact source dictionary and unresolved comparisons.}
To prevent a collision with the final matrix scheme $\mathcal M$
and the fibre-product group $\mathcal M_D$ of B.13, write
$\mathscr N$ for the source $\widetilde N_{1,1}$ and
$\mathscr P$ for the source $\widetilde{\mathcal M}_{0,3}$.
The source begins with
$1\to\langle\omega\rangle\to\mathscr N\to\mathscr P\to1$
at C3:573(1). The formal groups $G,Z,\gamma$ have been constructed
above; identifying them with these named geometric groups is
an additional comparison. The remaining assertions are the following.

\begin{enumerate}
\item[(R1)] \textbf{C3:573--576(3)--(5),(13),582,622(4 bis).}
For the source-defined action
$\iota:\mathscr P\to\Aut(\overline C)$, prove
\[
       \iota:\mathscr P\xrightarrow{\sim}\overline{\mathcal A},
       \qquad \mathscr N\xrightarrow{\sim}G
       \quad\text{over }C\text{ and the central }\omega.
 \tag{B.16.R1}
\]
Our (B.16.4) and (B.16.6) give the unique peripheral lift and
the formal central extension once this geometric comparison
is supplied. They do not identify an earlier geometric
$\pi_1$ construction merely by giving its abstract automorphism
candidate the same name. Nor is $\gamma$ thereby identified
with the absolute Galois group. The further GT and Galois
questions are D.7.

\item[(R2)] \textbf{C3:615--618(132)--(143),623(8)--(12).}
The first additional unquotiented lifting question is exactly
whether the proved section $b_\gamma$ in (B.16.12) admits a
continuous homomorphic lift
\[
 \widetilde b_\gamma:\gamma\longrightarrow N_s^*,
 \qquad\widetilde b_\gamma\bmod\omega=b_\gamma.
 \tag{B.16.R2}
\]
Equivalently, the explicitly defined central extension
\[
 E_s=\gamma\times_{N_s^*/\omega}N_s^*,
 \qquad1\to\langle\omega\rangle\to E_s\to\gamma\to1
 \tag{B.16.R2a}
\]
must split. The homomorphism $b$ we proved on $Z^!$ does not
itself descend to $\gamma$, since $b(e^2)=\omega$.
The manuscript obtains a lift from a chosen source section
$\mathscr N(0)\simeq\gamma$ of $Z^!\to\gamma$.
That section and its agreement with the earlier geometric
construction have not been constructed here.

The additional source objects consequently have the following
precise conditional dictionary, after choosing that section
and making (R1):
\[
\begin{array}{c|l}
\text{source object and scan}&\text{defining subgroup or intersection}\\ \hline
\mathscr N^!\quad616(\mathrm{NB}),623(8)&\ker\kappa=G^!\\
\mathscr N_0^!\quad616(135),623(8)&\mathscr N(0)P_0\\
\mathscr N(1/2)\quad617(139),623(9)&N_s^*\cap\mathscr N_0^!\\
\mathscr N'(-\bar j)\quad618(142),624(17)&N_r^*\cap\mathscr N_0^!\\
\gamma'\quad618&\ker(\epsilon_3:\gamma\to\{\pm1\})
\end{array}
 \tag{B.16.R2b}
\]
The remaining assertions are that the third row maps
isomorphically onto $\gamma$, that the fourth row maps
isomorphically onto $\gamma'$, and that these are the particular
geometric sections at $1/2$ and $-\bar j$ in C3:627(25).
These are specific section comparisons, not consequences of
the matrix-scheme section calculations in B.5.

\item[(R3)] \textbf{C3:619--621(144)--(149),625--626(21)--(24).}
The other proved quotient section defines the exact central
extension
\[
 E_r=\gamma\times_{N_r^*/\omega}N_r^*,
 \qquad1\to\langle\omega\rangle\to E_r\to\gamma\to1.
 \tag{B.16.R3}
\]
The remaining comparison identifies this particular extension
with the source $\widetilde\gamma_{r,s}$ and
$\mathcal N_r^?$ and identifies its class with the pullback
of the finite cocycle on C3:596--610. D.8 proves that finite
cocycle and D.12 computes the stated arithmetic cup product;
neither alone identifies the present profinite extension.
The source's claims about its nontriviality and about its
restriction to named geometric sections require this comparison.

The diagram C3:628(26), folio~617 verso (parallel page~627), introduces different
objects again. Write $\mathscr E_{1,1,\Q}$ for the stack of
elliptic curves with origin and $\mathscr E'_{1,1,\Q}$ for its
rigidification by the central involution $[-1]$. These are the
source's $M_{1,1},M'_{1,1}$; neither is the affine group scheme
$\mathcal M$ of Part~B. Concretely, rigidification stackifies the
same objects with isomorphisms divided by this central $C_2$.
Its fibres are therefore locally nonempty groupoids with central
automorphism group $C_2=\mu_2$ over $\Q$, which is the gerbe
meaning required in that diagram. The $j$-invariant factors
through this map to the coarse $j$-line. A specified object of
the rigidification is additional data beyond its $j$-coordinate.
The question in scan~628 about existence of a curve of invariant
zero over $\Q$ has the explicit answer
\[
 E:y^2=x^3+1,\qquad c_4=0,\quad c_6=-864,\quad
 \Delta=-432\ne0,\quad j=c_4^3/\Delta=0.
 \tag{B.16.19}
\]
The point at infinity supplies the origin and the nonzero
discriminant proves smoothness. The invariant identity
$c_4^3-c_6^2=1728\Delta$ and the displayed values are checked
in \texttt{D\_motives.py}. Over an algebraic closure the
automorphism $(x,y)\mapsto(\zeta_3x,y)$ survives rigidification,
so that object is not identified with a scheme point merely
by dividing out $[-1]$. The source itself distinguishes a lift
of a rigidified object from a lift of its coarse invariant in
C3:629, folio~618 (parallel page~628). Both photographs were
inspected. The exact geometric question in (R3)
is the lift of the particular object or section specified by
the source, not the existence of an elliptic curve with that
coarse invariant. No integral extension of this stack comparison
at residue characteristic two is asserted here.

The further lift laws and filtered groups on C3:583--595 and
635--650 are separately recorded in B.18.
\end{enumerate}
The ``Open'' label applies to (R1)--(R3), as unfinished assertions
of this report. It does not label the established Herfort--Ribes
theorems or our normalized conjugator calculation as open.
